% ==================
% 5) FROM HAIRCUT TO LENDING LIMITS
% ==================
\section{From haircut to lending limits}\label{sec:policy}
The four outputs of Section~\ref{sec:overview} follow in closed form, in one pass with no circular dependency and no iteration. The liquidation threshold and the borrow limit come from the aggregation engine of Section~\ref{sec:normalisation}, through the recipe of Eq.~\eqref{eq:recipe}. The two caps come from attested capacity, below. The policy layer adds five elements to the recipe. Three act on the operating range: the static liquidation buffer $\nu^{(a)}$ (Section~\ref{sec:policy-notch}), the due-diligence tier multiplier $\kappa^{(a)}$ (Section~\ref{sec:policy-tiers}), and the qualifying-insurance credit against the named operational layers (Section~\ref{sec:policy-insurance}). Two act on quantity and on state: the borrow and supply caps with their launch ramp (Section~\ref{sec:policy-caps}), and the valuation-native guardrails with their hard overrides (Section~\ref{sec:policy-guardrails}).

The clearing discount $\delta_{\mathrm{clear}}^{(a)}$ takes its single defining equation from Eq.~\eqref{eq:m-scaling}, evaluated on the clearing window $\mathrm{LH}^{(a)}$ that the three-source hierarchy of Section~\ref{sec:theory:setting} fixes. It is an additive haircut component \emph{inside} $H$, above the liquidation threshold. A window beyond ten business days makes it a charge. Contractually enforceable terms, attested or documented, that enforce a shorter window make it a credit. Estimated legs and FRTB fallback windows never earn one.

Caps follow from the issuer's attested redemption capacity per dealing window $R_{\mathrm{cap}}^{(a)}$ (Section~\ref{sec:policy-caps}). Capacity is the named exception to the floors of Section~\ref{sec:agg-missing}. Documented capacity evidence is a listing precondition rather than an input that takes a floor. The offering-memorandum proxy supports the first ramp step. Further exposure requires a signed attestation. An uncured lapse or refusal sets $R_{\mathrm{cap}}^{(a)}=0$, which sets both caps to zero.

The split of the price-risk charge into $\mathrm{IM}^{(a)}$ and $\delta_{\mathrm{clear}}^{(a)}$ leaves the reserve sufficient at the point of liquidation and costs nothing numerically. The argument for that sufficiency and the three reasons for the split are stated once, in Section~\ref{sec:theory}.

\subsection{Static liquidation buffer}\label{sec:policy-notch}
The liquidation buffer $\nu^{(a)}$ is fixed rather than recomputed from live conditions. It is a per-listing constant, the gap between $\theta^{(a)}$ and $\ell_{\max}^{(a)}$, inside the 4--8pp range of Section~\ref{sec:theory:limits}. The reference check of Eq.~\eqref{eq:notch-buffer} tests the adopted constant and does not recompute it from live data. A listing whose implied minimum exceeds its adopted buffer launches with a tighter ceiling for new borrows only, through the guardrails of Section~\ref{sec:policy-guardrails}, rather than with a wider buffer. Assets with sparse published valuations, or with a documented exposure that a single adverse valuation could realise as a jump, are the listings that need the tighter ceiling. A state-dependent gap is procyclical: it shrinks the buffer between live positions and the liquidation threshold exactly when published valuations fall and the redemption process is slowest, pushing performing positions into liquidation at the worst point of the cycle. Stress response instead uses those same guardrails, which can only tighten. None of them retroactively liquidates a performing position.

\subsection{Due-diligence tier multipliers}\label{sec:policy-tiers}
Eligibility and final limits pass through the due-diligence ranking framework (eligibility failures first, then the weighted ordinal score over seven dimensions, with the full rubric in the practitioner's scorecard). The dimensions are legal structure, redemption firmness, regulatory status, transparency, operational security, custody, and governance. An asset with a bridge or cross-chain dependence sits outside the scope of this framework and is not evaluated under it. That covers a bridged or wrapped token, and mint/burn, transfer agency, or a valuation oracle depending on cross-chain messaging. No score rescues an eligibility failure. For assets that pass, the tier multiplier $\kappa^{(a)} \in (0,1]$ scales the type's threshold, $\theta^{(a)} = \kappa^{(a)}\bigl(1-H^{(a)}\bigr)$, and the borrow LTV follows as $\ell_{\max}^{(a)} = \theta^{(a)} - \nu^{(a)}$, with the liquidation buffer $\nu^{(a)}$ absolute rather than tier-scaled. The over-collateralisation cushion at the point of liquidation is therefore $1-\theta^{(a)} \ge H^{(a)}$, strict below the Prime tier.

The tier structure, the thresholds, and the multiplier per tier are the practitioner's own elections, fixed in their scorecard (free parameter, Appendix~\ref{sec:free-parameters}). An illustrative structure has three tiers and a reject decision. Prime is a score of at least 80 with a minimum of 4 on legal structure, redemption firmness, and transparency, nothing below 3 on any dimension, and $\kappa=1.00$. Standard is a score of at least 65 with nothing below 2 on any dimension and $\kappa=0.90$. Restricted is a score of at least 50 with $\kappa=0.75$ and a concentration penalty the practitioner elects (free parameter, Appendix~\ref{sec:free-parameters}). Reject is a score below 50, any eligibility failure, or a zero on any dimension.

The tier acts on the over-collateralisation cushion only, and the layer table $\Lambda^{(a)}$ is tier-invariant. A scorecard dimension drives $\kappa^{(a)}$ through the weighted score or it drives a named layer, never both: each fact is charged exactly once. Transparency and governance drive $\kappa^{(a)}$. The effective oracle tier sets $\lambda_{\mathrm{oracle}}$ and caps the transparency score.

\subsection{Qualifying-insurance credit}\label{sec:policy-insurance}
Genuine, enforceable insurance reduces the charge. It can reduce \emph{only} the named operational/process layers it demonstrably covers, never the market or liquidation reserve. For a qualifying policy covering layer $\lambda_j^{(a)}$,
\begin{equation}\label{eq:insurance-credit}
\lambda_j^{(a),\mathrm{eff}} = \max\!\Bigl(\,\phi_j\,\lambda_j^{(a)},\;\; \lambda_j^{(a)} - R_{\mathrm{ins},j}^{(a)}\Bigr),
\qquad
R_{\mathrm{ins},j}^{(a)} = \beta_j\,\lambda_j^{(a)}\,\min\!\Bigl(1,\tfrac{\mathrm{NetLimit}}{E_{\mathrm{ref}}^{(a)}}\Bigr),
\end{equation}
where $E_{\mathrm{ref}}^{(a)}=\max(\text{live exposure},\,S_{\max}^{(a)})$, $\mathrm{NetLimit}$ is the per-occurrence limit \emph{net} of deductibles, retentions, shared-limit allocation, and excluded perils, $\beta_j\in[0,1]$ is the prescribed recognition factor for the policy's evidenced quality, and $\phi_j$ is a prescribed floor. The recognition factors, floors, rating minimum, and aggregate cap named below are practitioner elections registered in Appendix~\ref{sec:free-parameters}. The credit is bounded as follows:
\begin{enumerate}[label=(\roman*),itemsep=1pt]
\item \textbf{Eligible layers only:} the manager/operational layer $\lambda_{\mathrm{ops}}$ (custody, crime, cyber, or operational cover) and the oracle-integrity layer $\lambda_{\mathrm{oracle}}$ (professional-indemnity or valuation-error cover). It \emph{never} reduces the price-risk charge $\mathrm{IM}+\delta_{\mathrm{clear}}$, the valuation-staleness layer (baseline or its dynamic overdue-publication escalation), or the gating, default, wrong-way, or leverage layers.
\item \textbf{Capped and floored:} the reduction never exceeds the layer ($\beta_j\le 1$) and never takes it below $\phi_j\lambda_j$, where the practitioner sets the floor fraction $\phi_j$. One policy offsets exactly one layer (no stacking). The total reduction across all layers is bounded after the per-layer credits by an aggregate cap the practitioner sets. Where that aggregate cap governs, the reduction is allocated to the worst layer first.
\item \textbf{Qualification (else $\beta_j=0$):} a regulated insurer rated at or above the minimum the practitioner sets, in-force with premium paid, scope matched to a \emph{modelled} loss, for-benefit-of-lenders or loss-absorbing seniority, adequate limit, exclusions reviewed to confirm the policy pays in the modelled scenario, and a claims-paying record. Affiliated/captive or balance-sheet ``self-insurance'' scores $\beta_j=0$. Insurance never cures an eligibility failure (single-key control, sanctions, fraud, unrecovered loss) and never brings a bridge-dependent asset into scope.
\item \textbf{Recognition tiers}, set by the practitioner: standard qualifying, strong (broad wording, clean claims), and cash-collateralised / AA-insurer / letter-of-credit-equivalent. Grades are capped at $1-\phi_j$, the largest recognition the floor admits.
\item \textbf{Lapse:} non-payment, cancellation, rating downgrade below the rating minimum, limit exhaustion, or a relevant-peril exclusion withdraws the reduction immediately for new business ($R_{\mathrm{ins}}\to0$). The reduction returns only when the policy again qualifies.
\end{enumerate}
Until the practitioner sets $\beta_j$, $\phi_j$, and the rating minimum, $\beta_j=0$ and no reduction is granted. The graduated insurance evidence also feeds the legal-structure and operational-security dimensions of the due-diligence scorecard, replacing the binary top-level insurance condition. To prevent a double benefit, a policy recognised here for a layer reduction ($R_{\mathrm{ins}}>0$) supports the levels awarded on those dimensions only as disclosure and enforceability evidence, and it does \emph{not} additionally raise the tier multiplier $\kappa$ for the same coverage.

\subsection{Caps and launch ramp}\label{sec:policy-caps}
Both caps are keyed to the attested issuer per-window redemption capacity $R_{\mathrm{cap}}^{(a)}$. The same number re-keys the SIMM concentration threshold $T$ inside the concentration risk factor $\mathrm{CR}^{(a)} = \max\{1, \sqrt{\mathrm{position}^{(a)}/R_{\mathrm{cap}}^{(a)}}\}$ of Section~\ref{sec:normalisation}:
\begin{equation}\label{eq:policy-caps}
B_{\max}^{(a)} \;\le\; R_{\mathrm{cap}}^{(a)},
\qquad
S_{\max}^{(a)} \;\le\; M\, R_{\mathrm{cap}}^{(a)},
\end{equation}
with the policy constant $M$, the supply-cap multiple, a free parameter of the practitioner (Appendix~\ref{sec:free-parameters}). At the illustrative election $M=8$, $S_{\max}^{(a)} \le 8\,R_{\mathrm{cap}}^{(a)}$ and the re-keyed concentration risk factor is bounded at $\mathrm{CR} \le \sqrt{M} \approx 2.83$. The borrow cap is further bounded by the committed settlement level of Eq.~\eqref{eq:queue}, drawn from the subscription documents (Section~\ref{sec:efficacy}). $B_{\max}^{(a)}$ bounds borrowed \emph{debt} against collateral $a$ as a pool-level parameter of the deployment. The deployed market configuration shows the active borrow and supply caps. The practitioner sets $M$ as a model constant. $R_{\mathrm{cap}}^{(a)}$ and $S_{\max}^{(a)}$ are in the collateral's published-valuation USD numéraire.

\paragraph{Per-issuer global aggregation.}\label{sec:policy-issuer}
Eligibility, tiering, and capacity are judged once per issuer rather than per deployment. Where the same issuer's token is listed in multiple deployments, a single global issuer cap governs the sum of per-deployment caps,
\begin{equation}\label{eq:issuer-global-cap}
\sum_{\text{deployments}} S_{\max}^{(a)} \;\le\; S_{\mathrm{global}}^{\text{(issuer)}} \;\le\; M \, R_{\mathrm{cap}},
\end{equation}
because $R_{\mathrm{cap}}$ is an issuer-level resource that does not multiply with the number of listings. A policy consequence triggered at issuer level applies across every deployment that references the issuer.

The liquidator redeems collateral at the published valuation only until the debt is repaid. One dealing window must therefore clear the largest position's \emph{debt} at the point of liquidation, $B_{\max}^{(a)}\le R_{\mathrm{cap}}^{(a)}$, and excess collateral above the repaid debt queues at the borrower's risk (Section~\ref{sec:efficacy}). $R_{\mathrm{cap}}^{(a)}$ enters here on the one-window basis of Section~\ref{sec:theory:capacity}. A capacity attested on another basis is normalised to the governing dealing window.

During launch, the active on-chain cap may sit below Eq.~\eqref{eq:policy-caps}. The controlled operations record holds the launch schedule. Restricted-tier assets take the concentration penalty of Section~\ref{sec:policy-tiers}.

\subsection{Valuation-native guardrails and hard overrides}\label{sec:policy-guardrails}
Guardrails read observations from the redemption process itself: valuation publication, gates, capacity, and ratings. Every guardrail is automatic and can only tighten:
\begin{description}[leftmargin=1.5em,style=nextline,itemsep=1pt]
\item[Valuation not published within the contractual SLA.] Pause new borrows ($\ell_{\max}\to0$ for new positions). Existing positions are untouched. Recurrent breaches escalate the staleness layer.
\item[Gate, suspension, or holdback announced by the issuer.] Pause the asset (no new borrows or collateral inflow) and start an on-chain cap reduction toward orderly full redemption. $\theta$ on live positions is never lowered (Section~\ref{sec:overview-emergency}). A disorderly or unresolved redemption gate bars relisting. A resolved gate remains subject to controlled revalidation.
\item[Coverage-ratio breach (aggregate borrows exceed one $R_{\mathrm{cap}}^{(a)}$ window).] Cut $B_{\max}^{(a)}$ and $S_{\max}^{(a)}$.
\item[Securitised-collateral rating falls below investment grade.] New borrows pause and caps step down immediately. \emph{Confirmed} loss of investment grade reassigns new business to CreditNonQ bucket~2 and starts automatic full redemption. That bucket's published SIMM risk weight (1{,}300 at v2.6) raises $H^{(a)}$ toward the eligibility cutoff of Section~\ref{sec:agg-bounded}.
\item[Proof-of-reserves attestation failure or valuation restatement.] The failure can block relisting or tighten the due-diligence score.
\end{description}
The numeric alert levels behind the guardrails above, and the escalation level each action belongs to, are specified in the operations companion (Section~\ref{sec:operations}). The stablecoin peg on the borrow leg carries no charge under the USD-cash borrow policy. Hard overrides take precedence over the model: when any of the conditions above is met, its action applies regardless of what Eq.~\eqref{eq:recipe} or the tier multipliers would otherwise permit. Restoration to model-implied parameters occurs only after the guardrail condition clears.

\subsection{Operating rules and escalation}\label{sec:operations}
Automatic actions respond to adverse evidence and can only tighten active limits or pause new exposure. They never lower the liquidation threshold on a live position. An active value can loosen only after the condition that tightened it has cleared. The operations companion specifies the alert levels, monitoring, retention, procedures, and approvals.
